%%
%% paper79_temperley_lieb_jones_de.tex
%% Autor: Michael Fuhrmann
%% Projekt: Specialist Mathematics Project
%% Thema: Die Temperley-Lieb-Algebra und das Jones-Polynom: Offene Probleme
%% Sprache: Deutsch
%% Erstellt: 2026-03-12
%% Letzte Änderung: 2026-03-12
%% Build: 164
%%
%% Beschreibung:
%%   Deutschsprachige Version von Paper 79. Die Temperley-Lieb-Algebra
%%   TL_n(δ), das Jones-Polynom V_L(t), Khovanov-Homologie als
%%   Kategorifizierung, offene Probleme: Jones-Unknoten-Vermutung,
%%   Volumen-Vermutung (Kashaev-Murakami-Murakami).
%%   Die Jones-Unknoten-Vermutung und die Volumen-Vermutung sind offen.
%%

\documentclass[12pt,a4paper]{amsart}

%% Pakete für Mathematik, Layout, Encoding und Links
\usepackage{amsmath,amssymb,amsthm,mathtools,enumitem,booktabs,array,hyperref}
\usepackage[utf8]{inputenc}
\usepackage[T1]{fontenc}
\usepackage[ngerman]{babel}
\usepackage{geometry}
\geometry{margin=2.5cm}

%% Theorem-Umgebungen (englisch)
\newtheorem{theorem}{Theorem}[section]
\newtheorem{lemma}[theorem]{Lemma}
\newtheorem{corollary}[theorem]{Korollar}
\newtheorem{proposition}[theorem]{Proposition}

%% Deutsche Theorem-Umgebungen
\newtheorem{satz}[theorem]{Satz}
\newtheorem{korollar}[theorem]{Korollar}
\newtheorem{vermutung}[theorem]{Vermutung}

\theoremstyle{definition}
\newtheorem{definition}[theorem]{Definition}
\newtheorem{definition_de}[theorem]{Definition}
\newtheorem{beispiel}[theorem]{Beispiel}

\theoremstyle{remark}
\newtheorem{bemerkung}[theorem]{Bemerkung}
\newtheorem{remark}[theorem]{Bemerkung}
\newtheorem{conjecture}[theorem]{Vermutung}

%% Mathematische Operatoren und Abkürzungen
\newcommand{\Z}{\mathbb{Z}}
\newcommand{\Q}{\mathbb{Q}}
\newcommand{\C}{\mathbb{C}}
\DeclareMathOperator{\tr}{tr}
\DeclareMathOperator{\Vol}{Vol}
\DeclareMathOperator{\Kh}{Kh}

%% Hyperref-Konfiguration
\hypersetup{
    colorlinks=true,
    linkcolor=blue,
    citecolor=blue,
    urlcolor=blue,
    pdftitle={Die Temperley-Lieb-Algebra und das Jones-Polynom: Offene Probleme},
    pdfauthor={Michael Fuhrmann}
}

%% -------------------------------------------------------------------------
%% Abstract VOR \maketitle
%% -------------------------------------------------------------------------
\begin{abstract}
Die Temperley--Lieb-Algebra $\mathrm{TL}_n(\delta)$, eingeführt von Lieb und
Temperley im Rahmen der statistischen Mechanik~\cite{TemperleyLieb1971},
liefert den algebraischen Rahmen für das Jones-Polynom $V_L(t)$, eine
wegweisende Knoteninvariante, entdeckt von Vaughan Jones 1984~\cite{Jones1985}.
Wir präsentieren den diagrammatischen Kalkül von $\mathrm{TL}_n(\delta)$,
die Markov-Spur und ihre Rolle bei der Definition von $V_L(t)$, die Skeин-
Relation und die Haupteigenschaften des Jones-Polynoms.

Wir diskutieren dann die Khovanov-Homologie~\cite{Khovanov2000} als
\emph{Kategorifizierung} des Jones-Polynoms: eine bigraduierte Homologietheorie,
deren graduierte Euler-Charakteristik $V_L(t)$ wiedergibt.

Die zentralen offenen Probleme sind:
\begin{itemize}
\item Die \emph{Jones-Unknoten-Vermutung}: Impliziert $V_L(t) = 1$, dass $L$
      der triviale Knoten ist?
\item Die \emph{Volumen-Vermutung} von Kashaev--Murakami--Murakami: Ist die
      exponentielle Wachstumsrate des gefärbten Jones-Polynoms gleich dem
      hyperbolischen Volumen des Knotenkomplements?
\item Ob Khovanov-Homologie den trivialen Knoten erkennt (bekannt:
      ja, Kronheimer--Mrowka~\cite{KronheimerMrowka2011}).
\end{itemize}
\end{abstract}

\title{Die Temperley--Lieb-Algebra und das Jones-Polynom:\\
Offene Probleme}
\author{Michael Fuhrmann}
\address{Specialist Mathematics Project, Build 165}
\email{specialist-maths@project.local}
\date{2026-03-12}

\maketitle

%% -------------------------------------------------------------------------
\section{Einleitung}
%% -------------------------------------------------------------------------

Im Jahr 1984 entdeckte Vaughan Jones bei der Arbeit an von-Neumann-Algebren
eine neue Polynom-Invariante von Knoten und Verschlingungen~\cite{Jones1985}.
Das Jones-Polynom $V_L(t)$ ist ein Laurent-Polynom in $t^{1/2}$, das eine
topologische Invariante der orientierten Verschlingung $L$ ist. Seine
Entdeckung revolutionierte die Knotentheorie, brachte Jones die Fields-Medaille
(1990) ein und eröffnete neue Verbindungen zu Quantengruppen, statistischer
Mechanik, konformer Feldtheorie und 3-dimensionaler Quantentopologie.

Das algebraische Fundament des Jones-Polynoms ist die
\emph{Temperley--Lieb-Algebra} $\mathrm{TL}_n(\delta)$, eine assoziative
Algebra mit Erzeugern $\{e_1, \ldots, e_{n-1}\}$, die eine schöne
diagrammatische Realisierung durch „Tangle-Diagramme" besitzt. Das Jones-Polynom
wird durch eine Markov-Spur auf $\mathrm{TL}_n(\delta)$ konstruiert,
die sich als invariant unter Zopfisotopie und Markov-Bewegungen herausstellt.

\begin{vermutung}[Jones-Unknoten-Vermutung]
\label{verm:Jones-unknot}
Wenn $V_L(t) = 1$, dann ist $L$ der triviale Knoten.
\end{vermutung}

Diese scheinbar einfache Frage — ob das Jones-Polynom den trivialen Knoten
erkennt — ist nach 40 Jahren noch ungelöst.

%% -------------------------------------------------------------------------
\section{Die Temperley-Lieb-Algebra}
%% -------------------------------------------------------------------------

\subsection{Definition und Relationen}

\begin{definition_de}[Temperley-Lieb-Algebra]
Sei $\delta \in \mathbb{C}$. Die \emph{Temperley--Lieb-Algebra}
$\mathrm{TL}_n(\delta)$ ist die unitäre assoziative $\mathbb{C}$-Algebra
mit Erzeugern $\{e_1, e_2, \ldots, e_{n-1}\}$ und Relationen:
\begin{enumerate}[label=\roman*)]
\item $e_i^2 = \delta \cdot e_i$ \hfill (Idempotent bis auf Skalar);
\item $e_i e_{i\pm 1} e_i = e_i$ \hfill (Temperley-Lieb-Relation);
\item $e_i e_j = e_j e_i$ für $|i-j| \geq 2$ \hfill (entfernte Kommutativität).
\end{enumerate}
\end{definition_de}

\begin{satz}
$\dim_\mathbb{C} \mathrm{TL}_n(\delta) = C_n = \frac{1}{n+1}\binom{2n}{n}$
(Catalan-Zahl) für generisches $\delta$.
\end{satz}

\begin{proof}
Eine Basis bilden die $(n,n)$-Temperley-Lieb-Diagramme: Isotopie\-klassen
planarer Diagramme mit $n$ Punkten oben und $n$ Punkten unten, verbunden
durch nicht-kreuzende Bögen. Die Multiplikation ist durch vertikale
Konkatenation (mit Auswertung von Schleifen zu $\delta$) gegeben.
Die Anzahl solcher Diagramme ist die Catalan-Zahl $C_n$.
\end{proof}

\subsection{Diagrammatischer Kalkül}

Die Temperley--Lieb-Algebra hat eine natürliche graphische Darstellung:
\begin{itemize}
\item Die Identität ist durch $n$ vertikale Stränge repräsentiert.
\item Der Erzeuger $e_i$ ist durch ein Diagramm repräsentiert, bei dem
      Stränge $i$ und $i+1$ oben durch eine Tasse und unten durch einen
      Deckel verbunden sind.
\item Multiplikation entspricht vertikalem Stapeln.
\item Eine geschlossene Schleife wird zu $\delta$ ausgewertet.
\end{itemize}

\subsection{Einbettung in die Hecke-Algebra}

\begin{definition_de}[Iwahori-Hecke-Algebra]
Die \emph{Hecke-Algebra} $H_n(q)$ ist der Quotient der Gruppenalgebra
$\mathbb{C}[B_n]$ (Zopfgruppen-Algebra) durch die Relationen
$(T_i - q)(T_i + 1) = 0$.
\end{definition_de}

\begin{satz}[Jones 1987~\cite{Jones1987}]
Es gibt einen surjektiven Algebrahomomorphismus
$\rho \colon H_n(q) \to \mathrm{TL}_n(\delta)$
mit $\delta = q + q^{-1}$, gegeben durch $T_i \mapsto q e_i - 1$.
\end{satz}

%% -------------------------------------------------------------------------
\section{Das Jones-Polynom}
%% -------------------------------------------------------------------------

\subsection{Konstruktion via Markov-Spur}

\begin{definition_de}[Markov-Spur]
Eine \emph{Markov-Spur} auf $\bigcup_n \mathrm{TL}_n(\delta)$ ist eine
Familie linearer Abbildungen $\tau_n \colon \mathrm{TL}_n(\delta) \to \mathbb{C}$
mit:
\begin{enumerate}[label=\roman*)]
\item $\tau_n(ab) = \tau_n(ba)$ (Spureigenschaft);
\item $\tau_{n+1}(a \cdot e_n) = z \cdot \tau_n(a)$ für alle
      $a \in \mathrm{TL}_n(\delta)$ (Markov-Eigenschaft).
\end{enumerate}
\end{definition_de}

\begin{satz}[Jones 1985~\cite{Jones1985}]
Es existiert eine eindeutige Markov-Spur auf $\bigcup_n \mathrm{TL}_n(\delta)$
(bis auf Normierung), gegeben durch die „Schließung" von Diagrammen.
\end{satz}

\begin{definition_de}[Jones-Polynom]
Mit $\delta = -t^{1/2} - t^{-1/2}$ definieren wir:
\[
    V_L(t) = \left( -\frac{t^{1/2} + t^{-1/2}}{\delta} \right)^{n-1}
    \cdot \left( -t^{3/4} \right)^{w(\beta)} \cdot \tau_n(\rho(\beta)),
\]
wobei $\beta \in B_n$ ein Zopf mit Schließung $\hat{\beta} = L$ ist und
$w(\beta)$ die Verschlingungszahl (algebraische Länge) von $\beta$ ist.
Dies ist unabhängig vom Zopf-Repräsentanten (nach dem Satz von Markov).
\end{definition_de}

\subsection{Die Skein-Relation}

\begin{satz}[Jones-Skein-Relation~\cite{Jones1985}]
Das Jones-Polynom erfüllt die Skein-Relation:
\[
    t^{-1} V_{L_+}(t) - t \, V_{L_-}(t) = (t^{1/2} - t^{-1/2}) V_{L_0}(t),
\]
wobei $L_+$, $L_-$, $L_0$ drei Verschlingungen sind, die sich an einer
einzigen Kreuzung unterscheiden (positive, negative und geglättete Kreuzung).
\end{satz}

\begin{proof}
Dies folgt aus der Hecke-Relation $(T_i - q)(T_i + 1) = 0$, die
$T_i - T_i^{-1} = (q - q^{-1}) \cdot 1$ liefert. Übersetzt in
Verschlingungsdiagramme via der Markov-Spur ergibt dies die Skein-Relation.
\end{proof}

\subsection{Eigenschaften des Jones-Polynoms}

\begin{satz}[Grundlegende Eigenschaften]
Das Jones-Polynom $V_L(t)$ erfüllt:
\begin{enumerate}[label=\roman*)]
\item $V_{\text{trivial}}(t) = 1$;
\item $V_L(t) \in \Z[t^{1/2}, t^{-1/2}]$;
\item $V_{L_1 \sqcup L_2}(t) = (-t^{1/2} - t^{-1/2}) V_{L_1}(t) V_{L_2}(t)$;
\item $V_{\bar{L}}(t) = V_L(t^{-1})$ (Spiegelbild);
\item $V_{L_1 \# L_2}(t) = V_{L_1}(t) \cdot V_{L_2}(t)$.
\end{enumerate}
\end{satz}

\begin{beispiel}[Kleeblattknoten]
Für den rechtshändigen Kleeblattknoten $T(2,3)$:
\[
    V_{T(2,3)}(t) = -t^{-4} + t^{-3} + t^{-1}.
\]
Für den linkshändigen $T(2,-3)$ (Spiegelbild):
\[
    V_{T(2,-3)}(t) = -t^{4} + t^{3} + t.
\]
Da diese verschieden sind, ist der Kleeblattknoten chiral (nicht amphi-chiral).
\end{beispiel}

%% -------------------------------------------------------------------------
\section{Die Jones-Unknoten-Vermutung}
%% -------------------------------------------------------------------------

\begin{vermutung}[Jones-Unknoten-Vermutung~\cite{Jones1985}]
\label{verm:Jones}
Wenn $V_L(t) = 1$, dann ist $L$ der triviale Knoten.
\end{vermutung}

Dies ist das fundamentalste offene Problem über das Jones-Polynom.

\begin{bemerkung}
Die Vermutung ist bekanntermaßen gültig für:
\begin{itemize}
\item Alle Knoten bis zu 19 Kreuzungen (rechnerisch verifiziert).
\item Alternierende Knoten: Das Jones-Polynom eines nichttrivialen
      alternierenden Knotens kann nicht 1 sein (folgt aus den Tait-Vermutungen,
      bewiesen von Kauffman~\cite{Kauffman1987}, Murasugi~\cite{Murasugi1987},
      Thistlethwaite~\cite{Thistlethwaite1987}).
\item Positive Knoten (Tanaka~\cite{Tanaka1996}).
\end{itemize}
Kein Gegenbeispiel ist bekannt.
\end{bemerkung}

\begin{satz}[Kauffman-Klammer~\cite{Kauffman1987}]
Die \emph{Kauffman-Klammer} $\langle L \rangle \in \Z[A, A^{-1}]$ ist durch
die Skein-Relationen definiert:
\[
    \langle \bigcirc \rangle = 1, \quad
    \langle \bigcirc \sqcup D \rangle = (-A^2 - A^{-2})\langle D \rangle,
\]
\[
    \langle L_+ \rangle = A \langle L_0 \rangle + A^{-1} \langle L_\infty \rangle.
\]
Das Jones-Polynom ist dann
$V_L(t) = (-A)^{-3w(D)} \langle D \rangle \big|_{A = t^{-1/4}}$.
\end{satz}

%% -------------------------------------------------------------------------
\section{Das Gefärbte Jones-Polynom}
%% -------------------------------------------------------------------------

\begin{definition_de}[Gefärbtes Jones-Polynom]
Für eine irreduzible Darstellung $V_N$ von $\mathrm{SU}(2)$ der Dimension $N$
ist das \emph{$N$-gefärbte Jones-Polynom} $J_N(L; q)$ die Quanten-Invariante
von $L$, die durch Färben jeder Komponente mit $V_N$ entsteht.
Für $N=2$ ergibt sich das gewöhnliche Jones-Polynom $V_L(t)$ (bis auf
Normierung).
\end{definition_de}

%% -------------------------------------------------------------------------
\section{Die Volumen-Vermutung}
%% -------------------------------------------------------------------------

\begin{definition_de}[Kashaev-Invariante]
Für einen Knoten $K$ ist die \emph{Kashaev-Invariante} $\langle K \rangle_N$
bei der $N$-ten Einheitswurzel $q = e^{2\pi i/N}$ durch eine spezifische
$R$-Matrix-Konstruktion~\cite{Kashaev1997} definiert. Es gilt:
$\langle K \rangle_N = J_N(K; e^{2\pi i/N})$.
\end{definition_de}

\begin{vermutung}[Volumen-Vermutung, Kashaev 1997~\cite{Kashaev1997}, Murakami--Murakami 2001~\cite{MurakamiMurakami2001}]
\label{verm:Volumen}
Für einen hyperbolischen Knoten $K$ mit hyperbolischem Volumen
$\Vol(S^3 \setminus K)$ gilt:
\[
    \lim_{N \to \infty} \frac{2\pi \log |J_N(K; e^{2\pi i/N})|}{N}
    = \Vol(S^3 \setminus K).
\]
\end{vermutung}

\begin{beispiel}[Achtknoten]
Für den Achtknoten $4_1$ (den einfachsten hyperbolischen Knoten):
\[
    \Vol(S^3 \setminus 4_1) = 3v_3 \approx 2{,}0298,
\]
wobei $v_3 \approx 1{,}0149$ das Volumen des regulären idealen Tetraeders ist.
Die Volumen-Vermutung wurde für $4_1$ und viele andere Knoten numerisch
mit hoher Genauigkeit verifiziert.
\end{beispiel}

\begin{satz}[Murakami--Murakami~\cite{MurakamiMurakami2001}]
Für Torusknoten $T(2,2k+1)$ gilt:
\[
    \lim_{N \to \infty} \frac{2\pi \log |J_N(T(2,2k+1); e^{2\pi i/N})|}{N} = 0,
\]
konsistent mit der Volumen-Vermutung (Torusknoten sind nicht hyperbolisch;
ihr Komplement hat simpliziales Volumen null).
\end{satz}

\begin{bemerkung}
Die Volumen-Vermutung repräsentiert eine tiefe Verbindung zwischen
Quantentopologie und hyperbolischer Geometrie (Thurstons Geometrisierungsprogramm).
Sie ist für einige spezifische hyperbolische Knoten bewiesen, aber der
allgemeine Fall bleibt offen.
\end{bemerkung}

%% -------------------------------------------------------------------------
\section{Khovanov-Homologie}
%% -------------------------------------------------------------------------

\subsection{Kategorifizierung des Jones-Polynoms}

\begin{satz}[Khovanov 2000~\cite{Khovanov2000}]
\label{satz:Khovanov}
Es existiert eine bigraduierte Homologietheorie $\Kh^{i,j}(L)$,
kombinatorisch aus einem Diagramm von $L$ definiert, mit:
\begin{enumerate}[label=\roman*)]
\item $\Kh^{i,j}(L)$ ist für jedes $i,j \in \Z$ eine endlich erzeugte
      abelsche Gruppe und eine Invariante orientierter Verschlingungen;
\item Die graduierte Euler-Charakteristik liefert das Jones-Polynom:
      \[
          \sum_{i,j} (-1)^i q^j \dim \Kh^{i,j}(L)
          = (q + q^{-1}) V_L(-q^2);
      \]
\item Khovanov-Homologie ist streng stärker als das Jones-Polynom:
      Es gibt Paare von Verschlingungen mit gleichem Jones-Polynom
      aber verschiedener Khovanov-Homologie.
\end{enumerate}
\end{satz}

\subsection{Konstruktion}

Die Khovanov-Homologie wird wie folgt konstruiert:
\begin{enumerate}[label=\arabic*.]
\item Wähle ein Diagramm $D$ von $L$ mit $n$ Kreuzungen; bezeichne sie $1, \ldots, n$.
\item Für jede Teilmenge $S \subseteq \{1,\ldots,n\}$ bilde die \emph{$S$-Auflösung}
      $D_S$ (0-Glättung bei Kreuzungen außerhalb $S$, 1-Glättung bei Kreuzungen in $S$);
      dies ergibt eine Menge von Kreisen.
\item Weise jeder Auflösung den Vektorraum $V^{\otimes c(D_S)}$ zu,
      wobei $V = \Z\{v_+, v_-\}$ und $c(D_S)$ die Anzahl der Kreise ist.
\item Bilde den \emph{Khovanov-Komplex} $C^*(D)$ mit Differential
      $d \colon C^i \to C^{i+1}$ durch Multiplikations- und Komultiplikationsabbildungen
      (aus der Frobenius-Algebra-Struktur auf $V$).
\item $\Kh^*(D) = H^*(C^*(D))$ ist die Khovanov-Homologie.
\end{enumerate}

\begin{satz}[Bar-Natan 2002~\cite{BarNatan2002}]
Der Khovanov-Kettenkomplex $C^*(D)$ ist bis auf Kettenhomotopie-Äquivalenz
wohldefiniert (unabhängig von Diagrammwahlen).
\end{satz}

\subsection{Khovanov-Homologie erkennt den trivialen Knoten}

\begin{satz}[Kronheimer--Mrowka 2011~\cite{KronheimerMrowka2011}]
\label{satz:KM}
Khovanov-Homologie erkennt den trivialen Knoten: Falls
$\Kh(L) \cong \Kh(\text{trivial})$, dann ist $L$ der triviale Knoten.
\end{satz}

\begin{proof}[Beweisidee]
Kronheimer und Mrowka verwenden Instanton-Floer-Homologie. Sie definieren
eine Variante $\Kh'(L)$ (reduzierte Khovanov-Homologie mit $\Z/2$-Koeffizienten)
und zeigen ihre Verbindung zur Instanton-Knoten-Homologie $I^\#(L)$ via
eine Spektralsequenz. Die Instanton-Homologie erkennt den trivialen Knoten
(da $\dim I^\#(K) = 1$ genau für triviale Knoten $K$), und die
Spektralsequenz kollabiert auf der $E_2$-Seite für den trivialen Knoten.
\end{proof}

Satz~\ref{satz:KM} zeigt, dass Khovanov-Homologie streng stärker als das
Jones-Polynom ist. Folgendes bleibt offen:

\begin{vermutung}[Jones-Polynom erkennt trivialen Knoten, Verm.~\ref{verm:Jones}]
Wenn $V_L(t) = 1$, dann ist $L$ der triviale Knoten.
\end{vermutung}

%% -------------------------------------------------------------------------
\section{Ungerade Khovanov-Homologie}
%% -------------------------------------------------------------------------

\begin{definition_de}[Ungerade Khovanov-Homologie~\cite{OzsvathRasmussenSzabo2013}]
Die \emph{ungerade Khovanov-Homologie} $\Kh_{\mathrm{odd}}(L)$ ist eine
Variante der Khovanov-Homologie, die eine andere Frobenius-Algebra-Struktur
verwendet (äußere Algebra statt symmetrischer Algebra auf $V$). Es gilt:
\begin{enumerate}[label=\roman*)]
\item $\Kh_{\mathrm{odd}}^{i,j}(L)$ ist eine Verschlingungsinvariante;
\item Über $\Q$: $\Kh_{\mathrm{odd}}(L; \Q) \cong \Kh(L; \Q)$;
\item Über $\Z$: $\Kh_{\mathrm{odd}}(L)$ und $\Kh(L)$ können verschieden
      sein (unterschiedliche $\Z/2$-Torsion).
\end{enumerate}
\end{definition_de}

%% -------------------------------------------------------------------------
\section{Die Rasmussen $s$-Invariante aus der Khovanov-Homologie}
%% -------------------------------------------------------------------------

\begin{satz}[Rasmussen 2010~\cite{Rasmussen2010}]
Aus dem Khovanov-Komplex lässt sich eine Konkordanzinvariante
$s(L) \in 2\Z$, die \emph{Rasmussen $s$-Invariante}, gewinnen mit:
\begin{enumerate}[label=\roman*)]
\item $s$ ist ein Homomorphismus auf der Konkordanzgruppe;
\item $|s(K)|/2 \leq g_4(K)$ (glatte 4-Kugel-Geschlechts-Schranke);
\item $s$ liefert einen rein kombinatorischen Beweis, dass Torusknoten
      $T(p,q)$ das Geschlecht $g_4 = (p-1)(q-1)/2$ haben, der zuvor
      nur mit Eichtheorie (Kronheimer--Mrowka) zugänglich war.
\end{enumerate}
\end{satz}

Dies ist eine der bemerkenswertesten Anwendungen der Khovanov-Homologie:
ein kombinatorischer Beweis eines Resultats, das zuvor nur durch das
schwere Maschinerie der Eichtheorie erreichbar war.

%% -------------------------------------------------------------------------
\section{Offene Probleme}
%% -------------------------------------------------------------------------

\begin{vermutung}[Jones-Unknoten-Vermutung, Verm.~\ref{verm:Jones}]
$V_L(t) = 1$ genau dann, wenn $L$ der triviale Knoten ist.
\end{vermutung}

\begin{vermutung}[Volumen-Vermutung, Verm.~\ref{verm:Volumen}]
Für einen hyperbolischen Knoten $K$:
\[
    \lim_{N \to \infty} \frac{2\pi \log |J_N(K; e^{2\pi i/N})|}{N}
    = \Vol(S^3 \setminus K).
\]
\end{vermutung}

\begin{vermutung}[Gefärbte Volumen-Vermutung]
\label{verm:vol-gefaerbt}
Für das $N$-gefärbte Jones-Polynom bei $e^{2\pi i /N}$ kodiert auch
die Phase die Chern--Simons-Invariante:
\[
    J_N(K; e^{2\pi i /N}) \sim C \cdot N^{3/2} \cdot \exp\left(
    \frac{N}{2\pi}\bigl(\Vol(S^3 \setminus K)
    + i \cdot \mathrm{CS}(S^3 \setminus K)\bigr) \right).
\]
\end{vermutung}

\begin{vermutung}[AJ-Vermutung, Garoufalidis~\cite{Garoufalidis2004}]
\label{verm:AJ}
Das Quanten-$A$-Polynom (die Rekurrenzrelation für das gefärbte Jones-Polynom
$\{J_N(K;q)\}$) stimmt nach Spezialisierung $q \to 1$ mit dem klassischen
$A$-Polynom von $K$ überein.
\end{vermutung}

\begin{vermutung}[Funktorialität der Khovanov-Homologie]
\label{verm:Kh-func}
Khovanov-Homologie kann zu einem vollständig erweiterten topologischen
Quantenfeldtheorie-Funktor (TQFT) im Sinne von Lurie gemacht werden.
\end{vermutung}

\begin{vermutung}[Ungerade vs.\ gerade Khovanov-Homologie]
Es gibt primäre Knoten $K$, für die $\Kh(K; \Z)$ und
$\Kh_{\mathrm{odd}}(K; \Z)$ verschiedene $\Z$-Ränge haben (nicht nur
verschiedene Torsion).
\end{vermutung}

\begin{bemerkung}
Die Jones-Unknoten-Vermutung~\ref{verm:Jones} und die Volumen-Vermutung~\ref{verm:Volumen}
sind die beiden zentralen offenen Probleme in diesem Gebiet. Die
Unknoten-Vermutung würde folgen, wenn man zeigen könnte, dass jede
„exotische" Lösung von $V_L = 1$ durch Khovanov-Homologie oder eine
andere Invariante ausgeschlossen wird. Die Volumen-Vermutung verbindet
Quantentopologie und hyperbolische Geometrie auf eine Weise, deren
vollständiger mathematischer Inhalt noch erforscht wird.
\end{bemerkung}

%% -------------------------------------------------------------------------
\section{Zusammenfassung}
%% -------------------------------------------------------------------------

Die Temperley--Lieb-Algebra und das Jones-Polynom stellen einen Zusammenfluss
von Quantenalgebra, diagrammatischer Topologie und Theorie niedrig-dimensionaler
Mannigfaltigkeiten dar. Das Jones-Polynom unterscheidet den Kleeblattknoten
von seinem Spiegelbild, erkennt viele Knoten, die das Alexander-Polynom nicht
unterscheiden kann, und löste eine ganze Industrie neuer Invarianten aus
(HOMFLY, Kauffman-Klammer, Quantengruppen, topologische Quantenfeldtheorien).
Dennoch bleibt die einfachste Frage — impliziert $V_L = 1$, dass $L$ trivial ist? —
vollständig offen.

Khovanov-Homologie als Kategorifizierung des Jones-Polynoms ist streng
mächtiger (sie erkennt den trivialen Knoten, nach Kronheimer--Mrowka) und
ihr $s$-Invariant organisiert Informationen über Knotenkonkordanz. Die
Volumen-Vermutung verbindet Quanteninvarianten und hyperbolische Geometrie
auf eine profunde Weise, die noch vollständig zu verstehen bleibt.

\begin{thebibliography}{99}

\bibitem{TemperleyLieb1971}
H.~N.~V.~Temperley, E.~H.~Lieb,
\textit{Relations between the `percolation' and `colouring' problem and
other graph-theoretical problems associated with regular planar lattices},
Proc.\ Roy.\ Soc.\ London Ser.\ A \textbf{322} (1971), 251--280.

\bibitem{Jones1985}
V.~F.~R.~Jones,
\textit{A polynomial invariant for knots via von Neumann algebras},
Bull.\ Amer.\ Math.\ Soc.\ \textbf{12} (1985), 103--111.

\bibitem{Jones1987}
V.~F.~R.~Jones,
\textit{Hecke algebra representations of braid groups and link polynomials},
Ann.\ of Math.\ \textbf{126} (1987), 335--388.

\bibitem{Kauffman1987}
L.~H.~Kauffman,
\textit{State models and the Jones polynomial},
Topology \textbf{26} (1987), 395--407.

\bibitem{Murasugi1987}
K.~Murasugi,
\textit{Jones polynomials and classical conjectures in knot theory},
Topology \textbf{26} (1987), 187--194.

\bibitem{Thistlethwaite1987}
M.~B.~Thistlethwaite,
\textit{A spanning tree expansion of the Jones polynomial},
Topology \textbf{26} (1987), 297--309.

\bibitem{Khovanov2000}
M.~Khovanov,
\textit{A categorification of the Jones polynomial},
Duke Math.\ J.\ \textbf{101} (2000), 359--426.

\bibitem{BarNatan2002}
D.~Bar-Natan,
\textit{On Khovanov's categorification of the Jones polynomial},
Algebr.\ Geom.\ Topol.\ \textbf{2} (2002), 337--370.

\bibitem{Rasmussen2010}
J.~Rasmussen,
\textit{Khovanov homology and the slice genus},
Invent.\ Math.\ \textbf{182} (2010), 419--447.

\bibitem{KronheimerMrowka2011}
P.~B.~Kronheimer, T.~S.~Mrowka,
\textit{Khovanov homology is an unknot-detector},
Publ.\ Math.\ Inst.\ Hautes Études Sci.\ \textbf{113} (2011), 97--208.

\bibitem{Kashaev1997}
R.~M.~Kashaev,
\textit{The hyperbolic volume of knots from the quantum dilogarithm},
Lett.\ Math.\ Phys.\ \textbf{39} (1997), 269--275.

\bibitem{MurakamiMurakami2001}
H.~Murakami, J.~Murakami,
\textit{The colored Jones polynomials and the simplicial volume of a knot},
Acta Math.\ \textbf{186} (2001), 85--104.

\bibitem{OzsvathRasmussenSzabo2013}
P.~Ozsváth, J.~Rasmussen, Z.~Szabó,
\textit{Odd Khovanov homology},
Algebr.\ Geom.\ Topol.\ \textbf{13} (2013), 1465--1488.

\bibitem{Garoufalidis2004}
S.~Garoufalidis,
\textit{On the characteristic and deformation varieties of a knot},
Geom.\ Topol.\ Monogr.\ \textbf{7} (2004), 291--309.

\bibitem{Tanaka1996}
T.~Tanaka,
\textit{Maximal Thurston-Bennequin numbers of 2-bridge links},
Topology Appl.\ \textbf{68} (1996), 1--15.

\end{thebibliography}

\end{document}
